% =================================================
% 文件: 虚拟化.tex
% 章节: 虚拟化技术
% 版本: v1.0
% =================================================

\documentclass[12pt,UTF8]{ctexbook}

\input{../common/page}

\xeCJKsetup{AutoFallBack=true}
\setCJKfamilyfont{hei}{SimHei}
\setCJKfamilyfont{kai}{KaiTi}

\usepackage{titletoc}
\titlecontents{chapter}[0pt]{\vspace{3mm}\bf\addvspace{2pt}\filright}{\contentspush{\thecontentslabel\hspace{0.8em}}}{}{\titlerule*[8pt]{.}\contentspage}

\usepackage[colorlinks,linkcolor=blue,citecolor=blue,urlcolor=blue]{hyperref}

\ctexset{
	part/name= {第,卷},
	part/number={\chinese{part}},
	chapter/name={第,章},
	chapter/number={\arabic{chapter}}
}

\usepackage{graphicx}
\graphicspath{{Images/}}

\usepackage[perpage]{footmisc}

\usepackage{enumitem}

\title{\heiti\zihao{0} 虚拟化技术入门指南}
\author{WangFei}
\date{\today}

\begin{document}

\maketitle
\tableofcontents

\frontmatter
\chapter{前言}

本指南面向初学者，详细介绍虚拟化技术的基础概念、类型、实现方式、主流产品等知识。

\mainmatter

\chapter{虚拟化基础概念}
\label{chap:virtualization_basics}

\section{什么是虚拟化}
\label{sec:what_is_virtualization}

虚拟化是将物理资源抽象为虚拟资源的技术，使得多个虚拟环境可以共享物理硬件。

虚拟化的优势：
\begin{itemize}
    \item \textbf{资源利用率}：提高硬件资源利用率
    \item \textbf{隔离性}：虚拟环境相互隔离
    \item \textbf{灵活性}：快速创建和销毁虚拟环境
    \item \textbf{可移植性}：虚拟环境可在不同硬件间迁移
\end{itemize}

\section{虚拟化类型}
\label{sec:virtualization_types}

常见的虚拟化类型：
\begin{table}[htbp]
    \centering
    \caption{虚拟化类型}
    \begin{tabular}{|c|p{8cm}|}
        \hline
        \textbf{类型} & \textbf{说明} \\
        \hline
        完全虚拟化 & 虚拟机运行完整操作系统，无需修改 \\
        \hline
        半虚拟化 & 操作系统需要修改以配合虚拟化 \\
        \hline
        硬件辅助虚拟化 & 利用硬件特性加速虚拟化 \\
        \hline
        容器虚拟化 & 共享内核的轻量级虚拟化 \\
        \hline
        桌面虚拟化 & 将桌面环境虚拟化 \\
        \hline
    \end{tabular}
\end{table}

\section{虚拟化架构}
\label{sec:virtualization_architecture}

虚拟化架构主要有两种：

\begin{table}[htbp]
    \centering
    \caption{虚拟化架构}
    \begin{tabular}{|c|p{8cm}|}
        \hline
        \textbf{架构} & \textbf{说明} \\
        \hline
        Type 1 (裸金属) & 虚拟机监控器直接运行在硬件上 \\
        \hline
        Type 2 (宿主) & 虚拟机监控器运行在宿主操作系统上 \\
        \hline
    \end{tabular}
\end{table}

\chapter{KVM 虚拟化}
\label{chap:kvm}

\section{KVM 基础}
\label{sec:kvm_basics}

KVM（Kernel-based Virtual Machine）是 Linux 内核内置的虚拟化模块。

KVM 的特点：
\begin{itemize}
    \item 基于硬件辅助虚拟化
    \item 性能接近原生
    \item 支持多种操作系统
    \item 与 QEMU 配合使用
\end{itemize}

\section{安装 KVM}
\label{sec:install_kvm}

在 CentOS Stream 9 上安装 KVM：

\begin{verbatim}
dnf install -y qemu-kvm libvirt libvirt-daemon-kvm \
    virt-install virt-manager bridge-utils
systemctl enable --now libvirtd
\end{verbatim}

验证安装：
\begin{verbatim}
lsmod | grep kvm
virsh list --all
\end{verbatim}

\section{创建虚拟机}
\label{sec:create_vm}

使用 virt-install 创建虚拟机：

\begin{verbatim}
virt-install \
  --name centos9 \
  --ram 2048 \
  --vcpus 2 \
  --disk path=/var/lib/libvirt/images/centos9.qcow2,size=20 \
  --os-type linux \
  --os-variant centos9 \
  --network bridge=virbr0 \
  --graphics none \
  --console pty,target_type=serial \
  --location 'http://mirror.centos.org/centos-stream/9-stream/BaseOS/x86_64/os/' \
  --extra-args 'console=ttyS0,115200n8'
\end{verbatim}

管理虚拟机：
\begin{verbatim}
virsh start centos9
virsh shutdown centos9
virsh reboot centos9
virsh console centos9
\end{verbatim}

\section{网络配置}
\label{sec:network_config}

创建桥接网络：
\begin{verbatim}
nmcli connection add type bridge con-name br0 ifname br0
nmcli connection modify br0 ipv4.addresses 192.168.1.100/24
nmcli connection modify br0 ipv4.gateway 192.168.1.1
nmcli connection modify br0 ipv4.dns 192.168.1.1
nmcli connection modify br0 ipv4.method manual
nmcli connection add type bridge-slave con-name br0-port1 ifname eth0 master br0
nmcli connection up br0
\end{verbatim}

\chapter{QEMU 虚拟化}
\label{chap:qemu}

\section{QEMU 基础}
\label{sec:qemu_basics}

QEMU 是一个通用的开源虚拟机监控器。

QEMU 的特点：
\begin{itemize}
    \item 纯软件实现，无需硬件支持
    \item 支持多种架构
    \item 可与 KVM 配合使用
    \item 支持多种设备模拟
\end{itemize}

\section{QEMU 命令行}
\label{sec:qemu_command}

使用 QEMU 启动虚拟机：
\begin{verbatim}
qemu-system-x86_64 \
  -m 2048 \
  -smp 2 \
  -hda centos9.qcow2 \
  -cdrom CentOS-Stream-9-x86_64-dvd1.iso \
  -boot d \
  -net nic,model=virtio \
  -net user,hostfwd=tcp::2222-:22 \
  -vnc :0
\end{verbatim}

创建磁盘镜像：
\begin{verbatim}
qemu-img create -f qcow2 centos9.qcow2 20G
qemu-img resize centos9.qcow2 +10G
qemu-img info centos9.qcow2
\end{verbatim}

\chapter{LXC 容器}
\label{chap:lxc}

\section{LXC 基础}
\label{sec:lxc_basics}

LXC（Linux Containers）是 Linux 原生的容器化技术。

LXC 的特点：
\begin{itemize}
    \item 共享内核
    \item 轻量级
    \item 快速启动
    \item 资源隔离
\end{itemize}

\section{安装 LXC}
\label{sec:install_lxc}

在 CentOS Stream 9 上安装 LXC：

\begin{verbatim}
dnf install -y lxc lxc-templates
systemctl enable --now lxc
\end{verbatim}

创建容器：
\begin{verbatim}
lxc-create -n mycontainer -t centos
lxc-start -n mycontainer
lxc-attach -n mycontainer
\end{verbatim}

管理容器：
\begin{verbatim}
lxc-list
lxc-stop -n mycontainer
lxc-destroy -n mycontainer
\end{verbatim}

\chapter{Docker 容器}
\label{chap:docker}

\section{Docker 基础}
\label{sec:docker_basics}

Docker 是最流行的容器化平台。

Docker 的组件：
\begin{table}[htbp]
    \centering
    \caption{Docker 组件}
    \begin{tabular}{|c|p{8cm}|}
        \hline
        \textbf{组件} & \textbf{说明} \\
        \hline
        Docker Engine & 核心运行时 \\
        \hline
        Docker Hub & 镜像仓库 \\
        \hline
        Dockerfile & 镜像构建文件 \\
        \hline
        Docker Compose & 多容器编排 \\
        \hline
    \end{tabular}
\end{table}

\section{安装 Docker}
\label{sec:install_docker}

在 CentOS Stream 9 上安装 Docker：

\begin{verbatim}
dnf install -y dnf-utils
dnf config-manager --add-repo https://download.docker.com/linux/centos/docker-ce.repo
dnf install -y docker-ce docker-ce-cli containerd.io
systemctl enable --now docker
usermod -aG docker $USER
\end{verbatim}

验证安装：
\begin{verbatim}
docker run hello-world
\end{verbatim}

\section{Docker 命令}
\label{sec:docker_commands}

常用 Docker 命令：
\begin{verbatim}
docker images
docker ps -a
docker run -d -p 80:80 nginx
docker exec -it <container-id> bash
docker stop <container-id>
docker rm <container-id>
docker rmi <image-id>
\end{verbatim}

\section{Dockerfile}
\label{sec:dockerfile}

Dockerfile 示例：
\begin{verbatim}
FROM centos:stream9
RUN dnf install -y nginx
COPY index.html /usr/share/nginx/html/
EXPOSE 80
CMD ["nginx", "-g", "daemon off;"]
\end{verbatim}

构建镜像：
\begin{verbatim}
docker build -t mynginx .
\end{verbatim}

\section{Docker Compose}
\label{sec:docker_compose}

docker-compose.yml 示例：
\begin{verbatim}
version: '3.8'
services:
  web:
    image: nginx
    ports:
      - "80:80"
    volumes:
      - ./html:/usr/share/nginx/html
  db:
    image: mysql:8.0
    environment:
      MYSQL_ROOT_PASSWORD: password
    volumes:
      - db-data:/var/lib/mysql
volumes:
  db-data:
\end{verbatim}

启动服务：
\begin{verbatim}
docker-compose up -d
\end{verbatim}

\chapter{虚拟机管理}
\label{chap:vm_management}

\section{虚拟机迁移}
\label{sec:vm_migration}

虚拟机迁移分为两种：
\begin{itemize}
    \item \textbf{冷迁移}：关闭虚拟机后迁移
    \item \textbf{热迁移}：虚拟机运行时迁移
\end{itemize}

使用 virsh 进行热迁移：
\begin{verbatim}
virsh migrate --live centos9 qemu+ssh://remote-host/system
\end{verbatim}

\section{存储管理}
\label{sec:storage_management}

虚拟机存储类型：
\begin{table}[htbp]
    \centering
    \caption{存储类型}
    \begin{tabular}{|c|p{8cm}|}
        \hline
        \textbf{类型} & \textbf{说明} \\
        \hline
        qcow2 & QEMU 镜像格式，支持快照 \\
        \hline
        raw & 原始格式，性能最好 \\
        \hline
        LVM & 逻辑卷管理 \\
        \hline
        Ceph & 分布式存储 \\
        \hline
    \end{tabular}
\end{table}

创建快照：
\begin{verbatim}
virsh snapshot-create-as centos9 snapshot1
virsh snapshot-revert centos9 snapshot1
\end{verbatim}

\section{性能优化}
\label{sec:performance_optimization}

虚拟机性能优化：
\begin{itemize}
    \item 使用 virtio 驱动
    \item 启用 KVM 硬件加速
    \item 合理分配 CPU 和内存
    \item 使用 SSD 存储
    \item 优化网络配置
\end{itemize}

\chapter{云原生虚拟化}
\label{chap:cloud_native}

\section{容器编排}
\label{sec:container_orchestration}

容器编排工具：
\begin{table}[htbp]
    \centering
    \caption{容器编排工具}
    \begin{tabular}{|c|p{8cm}|}
        \hline
        \textbf{工具} & \textbf{说明} \\
        \hline
        Kubernetes & 最流行的容器编排平台 \\
        \hline
        Docker Swarm & Docker 原生编排工具 \\
        \hline
        Apache Mesos & 通用资源调度框架 \\
        \hline
    \end{tabular}
\end{table}

\section{Serverless}
\label{sec:serverless}

Serverless 架构：
\begin{itemize}
    \item 无需管理服务器
    \item 按需付费
    \item 自动扩展
    \item 事件驱动
\end{itemize}

常见的 Serverless 平台：
\begin{table}[htbp]
    \centering
    \caption{Serverless 平台}
    \begin{tabular}{|c|p{8cm}|}
        \hline
        \textbf{平台} & \textbf{说明} \\
        \hline
        AWS Lambda & AWS 无服务器计算服务 \\
        \hline
        Azure Functions & Azure 无服务器计算服务 \\
        \hline
        Google Cloud Functions & GCP 无服务器计算服务 \\
        \hline
        Knative & Kubernetes 原生 Serverless \\
        \hline
    \end{tabular}
\end{table}

\chapter{虚拟化安全}
\label{chap:virtualization_security}

\section{安全挑战}
\label{sec:security_challenges}

虚拟化环境的安全挑战：
\begin{itemize}
    \item 虚拟机逃逸
    \item 共享资源攻击
    \item 管理接口安全
    \item 镜像安全
\end{itemize}

\section{安全措施}
\label{sec:security_measures}

虚拟化安全措施：
\begin{itemize}
    \item 及时更新虚拟机监控器
    \item 隔离虚拟网络
    \item 安全配置虚拟机
    \item 使用加密存储
    \item 监控虚拟机行为
\end{itemize}

\chapter{虚拟化实践}
\label{chap:practice}

\section{搭建测试环境}
\label{sec:test_environment}

使用虚拟化搭建测试环境：
\begin{verbatim}
virt-install \
  --name test-server \
  --ram 4096 \
  --vcpus 4 \
  --disk path=/var/lib/libvirt/images/test-server.qcow2,size=50 \
  --os-type linux \
  --os-variant centos9 \
  --network bridge=br0 \
  --graphics vnc,listen=0.0.0.0
\end{verbatim}

\section{容器化应用}
\label{sec:containerize_app}

将应用容器化的步骤：
\begin{enumerate}
    \item 编写 Dockerfile
    \item 构建镜像
    \item 测试容器
    \item 推送镜像到仓库
    \item 部署到生产环境
\end{enumerate}

\backmatter

\end{document}